% !TeX root=../main.tex
\chapter{راهنمای کاربر}\label{Chap:UserGuide}

این فصل قابلیت‌های فعلی نوواژه و جریان استفاده از آن‌ها را از دید والد و کودک توضیح می‌دهد. برنامه امکان مرور واژگان فارسی همراه با تصویر و صوت، آزمون شنیداری و چند بازی کوتاه را فراهم می‌کند. والد می‌تواند نمایهٔ کودک و برخی محتوای شخصی را مدیریت کند. این فعالیت‌ها برای تمرین و ایجاد فرصت گفت‌وگو در کنار تعامل روزمره طراحی شده‌اند و جای ارزیابی یا راهنمایی متخصص را نمی‌گیرند.

\section{ورود به برنامه و انتخاب مسیر استفاده}

برنامه پس از اجرا داشبوردی را نشان می‌دهد که فعالیت‌های کودک از آن در دسترس‌اند. استفاده از آموزش واژه و بازی‌ها به ورود حساب وابسته نیست؛ بنابراین خانواده می‌تواند به‌صورت مهمان نیز فعالیت‌ها را اجرا کند. در این حالت، برنامه نمایهٔ کودک فعالی ندارد و سابقهٔ اختصاصی یک کودک برای مهمان ثبت نمی‌شود.

برای مدیریت کودک و دسترسی به تنظیمات، والد ابتدا وارد حساب می‌شود یا حساب تازه‌ای می‌سازد. پس از ورود، زبانهٔ تنظیمات برای والد باز می‌شود و نمایه‌های کودک و محتوای شخصی را در اختیار او قرار می‌دهد. با انتخاب نمایهٔ کودک، فعالیت‌های پشتیبانی‌شده به همان کودک نسبت داده می‌شوند. تفاوت این دو مسیر در مدیریت و نگهداری سابقه است: مهمان می‌تواند بازی کند، اما برای پیگیری پیشرفت کودک باید از حساب والد و نمایهٔ انتخاب‌شده استفاده شود.

\subsection{ساخت حساب و ورود}

صفحهٔ ورود نام کاربری و رمز عبور را دریافت می‌کند. برای ساخت حساب جدید، نام کاربری، نام نمایشی و رمز عبور وارد می‌شوند؛ نام نمایشی برای خطاب‌کردن والد در برنامه به‌کار می‌رود. پس از ورود یا ثبت‌نام موفق، برنامه به فضای اصلی والد بازمی‌گردد. اگر ورود انجام نشد، کامل‌بودن اطلاعات و دسترسی شبکه را بررسی و دوباره تلاش کنید.

\begin{figure}[H]
\centering
\begin{minipage}{0.47\textwidth}
\centering
\includegraphics[height=6.2cm]{\detokenize{figures/login_screen.png}}
\\[-1mm]\small (الف) صفحهٔ ورود
\end{minipage}
\hfill
\begin{minipage}{0.47\textwidth}
\centering
\includegraphics[height=6.2cm]{\detokenize{figures/signup_screen.png}}
\\[-1mm]\small (ب) صفحهٔ ثبت‌نام
\end{minipage}
\caption{ورود به حساب والد و ایجاد حساب جدید}
\label{Fig:UserAuth}
\end{figure}

\subsection{داشبورد والد و نمایهٔ کودک}

داشبورد فعالیت‌ها را در کارت‌های آزمون و بازی و در بخش آموزش دسته‌بندی‌شده ارائه می‌کند. با انتخاب یک دسته، مرور واژه‌های آن آغاز می‌شود؛ کارت‌های بازی نیز مستقیماً فعالیت مربوط را باز می‌کنند. در حالت مهمان، زبانهٔ ورود برای دسترسی به امکانات والد نمایش داده می‌شود.

والد از زبانهٔ تنظیمات می‌تواند نمایه‌های کودک را ایجاد یا انتخاب کند، واژه‌های شخصی را مدیریت کند، مدت جلسهٔ آموزش را تنظیم کند و راهنمای حالت کودک را ببیند. هنگام ایجاد نمایه، نام، سن و جنسیت دریافت می‌شوند. انتخاب کودک فعال برای ثبت پاسخ‌های آزمون و پیشرفت بازی‌های پشتیبانی‌شده اهمیت دارد؛ پیش از شروع فعالیت، والد باید مطمئن شود نمایهٔ مورد نظر انتخاب شده است. ورود دوباره به تنظیمات حساب، تأیید مجدد رمز عبور را می‌خواهد. با ترک زبانهٔ تنظیمات، دسترسی آن دوباره قفل می‌شود.

\begin{figure}[H]
\centering
\begin{minipage}{0.47\textwidth}
\centering
\includegraphics[height=6.3cm]{\detokenize{figures/main_screen_(mother_screen_ dashboard).png}}
\\[-1mm]\small (الف) داشبورد فعالیت‌ها
\end{minipage}
\hfill
\begin{minipage}{0.47\textwidth}
\centering
\includegraphics[height=6.3cm]{\detokenize{figures/mother_screen_setting.png}}
\\[-1mm]\small (ب) تنظیمات والد
\end{minipage}
\caption{داشبورد برنامه و تنظیمات حساب والد}
\label{Fig:UserDashboardSettings}
\end{figure}

\section{آموزش و آزمون واژگان}

\subsection{مرور واژه‌ها}

در بخش آموزش، والد یا کودک یک دستهٔ فعال را انتخاب می‌کند. هر واژه با تصویر و نام فارسی نمایش داده می‌شود و دکمهٔ پخش، تلفظ آن را ارائه می‌کند. پیکان‌های صفحه برای جابه‌جایی دستی میان واژه‌های قبلی و بعدی هستند. پس از پخش صوت، اگر کاربر اقدامی نکند، صفحه پس از چند ثانیه به واژهٔ بعدی می‌رود؛ این رفتار فقط به صفحهٔ آموزش مربوط است. زمان جلسه نیز در همین صفحه و بر پایهٔ مدت انتخاب‌شده در تنظیمات والد شمارش می‌شود و تایمر مشترک همهٔ بازی‌ها نیست.

در صورت انتخاب نمایهٔ کودک، شاخص‌های درون‌برنامه‌ای واژه و دسته می‌توانند در صفحهٔ آموزش نمایش داده شوند. این شاخص‌ها بازخوردی دربارهٔ استفاده از برنامه‌اند و به‌تنهایی نشان‌دهندهٔ تسلط کودک بر واژه در گفت‌وگوی روزمره نیستند. دسته‌های فعال از محتوای جاری دریافت می‌شوند؛ در نسخهٔ فعلی دستهٔ اعداد در دسترس است و دستهٔ قدیمی خودروهای ایرانی در فهرست فعال قرار ندارد.

\begin{figure}[H]
\centering
\begin{minipage}{0.47\textwidth}
\centering
\includegraphics[height=6.5cm]{\detokenize{figures/game_screen_(learning).png}}
\\[-1mm]\small (الف) نمایش واژه، تصویر و صوت
\end{minipage}
\hfill
\begin{minipage}{0.47\textwidth}
\centering
\includegraphics[height=6.5cm]{\detokenize{figures/quiz_screen.pmg.png}}
\\[-1mm]\small (ب) پرسش شنیداری و گزینه‌های تصویری
\end{minipage}
\caption{مرور واژه و نمونه‌ای از تعامل آزمون شنیداری}
\label{Fig:UserLearningQuiz}
\end{figure}

\subsection{آزمون شنیداری}

آزمون با انتخاب یک دسته یا حالت ترکیبی آغاز می‌شود و معمولاً پنج پرسش دارد. در هر پرسش کودک صوت واژه را می‌شنود و تصویر متناظر را از میان چهار گزینه انتخاب می‌کند. پاسخ درست یا نادرست بلافاصله بازخورد می‌گیرد و نوار پیشرفت جایگاه پرسش را در جلسه نشان می‌دهد. در صورت انتخاب نمایهٔ کودک، پاسخ‌ها و سابقهٔ آزمون برای همان کودک نگهداری می‌شوند؛ در حالت مهمان این سابقه به نمایه‌ای نسبت داده نمی‌شود. ممکن است واژه‌هایی که هنوز در شاخص برنامه آموخته‌شده ثبت نشده‌اند در اولویت پرسش قرار گیرند؛ این اولویت‌دهی محدود است و آزمون را به سنجش جامع یا سازوکار تطبیقی ارزیابی‌شده تبدیل نمی‌کند.

\section{بازی‌های کودک}

داشبورد پنج بازی کودک را در دسترس قرار می‌دهد: حافظه، مرتب‌سازی رنگ‌ها، تشخیص مورد متفاوت، تطبیق سایه و صورت و بدن. در حافظه، تشخیص مورد متفاوت و تطبیق سایه امکان انتخاب یک دسته یا حالت ترکیبی وجود دارد؛ مرتب‌سازی رنگ‌ها و صورت و بدن با محتوای درون خود بازی آغاز می‌شوند. در بازی‌هایی که ساختار مرحله‌ای دارند، ستاره‌ها معمولاً پیشرفت میان فعالیت‌های موفق همان سطح را نمایش می‌دهند. در «کدام متفاوت است؟» ستاره به درستی پاسخ در نخستین تلاش مربوط است و معنای نردبان سطحی ندارد.

\subsection{بازی حافظه}

بازی حافظه از انتخاب دسته یا مجموعهٔ ترکیبی از دسته‌ها آغاز می‌شود. کودک کارت‌های پوشیده را دو‌به‌دو باز می‌کند تا تصویرهای همسان را پیدا کند. جفت درست به‌عنوان تطبیق‌شده باقی می‌ماند؛ دو کارت ناهماهنگ پس از بازخورد کوتاه دوباره پوشانده می‌شوند. هر صفحه زمان محدود خود را دارد و تمام‌نشدن جفت‌ها پیش از پایان زمان، آن صفحه را ناموفق می‌کند.

بازی چهار بُعد دارد: صفحهٔ \lr{2×2} با دو جفت و ۳۰ ثانیه، \lr{2×3} با سه جفت و ۴۵ ثانیه، \lr{3×4} با شش جفت و ۶۰ ثانیه، و \lr{4×4} با هشت جفت و ۹۰ ثانیه. سه برد در هر بُعد، سه ستارهٔ آن بخش را کامل و بازی را به بُعد بعدی منتقل می‌کند. پس از سه برد در بُعد پایانی، بازی در همان بُعد می‌ماند و قابل تکرار است. با نمایهٔ کودک، پیشرفت بُعد برای ادامهٔ بعدی حفظ می‌شود؛ اجرای مهمان سابقهٔ پایدار کودک ایجاد نمی‌کند.

\begin{figure}[H]
\centering
\includegraphics[height=6.2cm]{\detokenize{figures/memory_game.png}}
\caption{نمونه‌ای از صفحهٔ بازی حافظه در بُعد ۴×۴}
\label{Fig:UserMemory}
\end{figure}

\subsection{بازی مرتب‌سازی رنگ‌ها}

در این بازی کودک بادکنک‌ها را با کشیدن و رهاکردن در سبد هم‌رنگ قرار می‌دهد. جای‌گذاری درست با بازخورد موفقیت تأیید می‌شود و خطا بازخورد موقت می‌گیرد تا کودک تلاش دیگری کند. سطح‌ها به‌ترتیب سه، چهار و پنج رنگ دارند و پس از سه صفحهٔ موفق، بازی به سطح بعد می‌رود. نوار ستاره پیشرفت سه صفحهٔ هر سطح را نشان می‌دهد. پس از تکمیل سطح سوم، بازی همان سطح را برای تکرار نگه می‌دارد؛ با انتخاب نمایه، پیشرفت سطح قابل بازیابی است.

\begin{figure}[H]
\centering
\includegraphics[height=6.2cm]{\detokenize{figures/color_game.png}}
\caption{مرتب‌سازی پنج رنگ در سبدهای هم‌رنگ در سطح سوم}
\label{Fig:UserColorSorting}
\end{figure}

\subsection{بازی تشخیص مورد متفاوت}

در هر دور چهار تصویر نمایش داده می‌شوند: سه مورد از یک دسته و یک مورد از دسته‌ای دیگر. کودک مورد متفاوت را انتخاب می‌کند. انتخاب نادرست بازخورد می‌گیرد و فرصت پاسخ دوباره در همان دور باقی می‌ماند؛ پس از انتخاب درست، دور بعد آغاز می‌شود. هر جلسه به‌طور پیش‌فرض پنج دور دارد و ستاره فقط برای پاسخی داده می‌شود که در نخستین تلاش درست باشد. برد نهایی زمانی اعلام می‌شود که تعداد ستاره‌ها بیش از نیمی از دورها باشد. این بازی نردبان سطح پایدار ندارد و ستاره‌ها نتیجهٔ همان جلسه‌اند.

\begin{figure}[H]
\centering
\includegraphics[height=6.2cm]{\detokenize{figures/odd_one_game.png}}
\caption{انتخاب تصویر متفاوت از سه تصویر هم‌دسته}
\label{Fig:UserOddOneOut}
\end{figure}

\subsection{بازی تطبیق سایه}

در این بازی تصویرهای سایه‌وار با تصویرهای گزینه تطبیق داده می‌شوند. کودک گزینه‌ای را انتخاب می‌کند که با هدف سایه‌ای سازگار است؛ پاسخ درست با نمایش بازخورد تأیید می‌شود و انتخاب نادرست بازخورد موقت می‌گیرد. انتخاب دسته یا حالت ترکیبی هنگام آغاز در دسترس است. سه سطح بازی به‌ترتیب یک هدف از چهار گزینه، دو هدف از چهار گزینه و سه هدف از شش گزینه دارند. پس از سه صفحهٔ موفق، مرحلهٔ بعد آغاز می‌شود و ستاره‌ها پیشرفت همان سطح را نشان می‌دهند. سطح سوم پایانی است و پس از تکمیل، بازی برای تکرار در همان سطح می‌ماند. پیشرفت سطح برای نمایهٔ کودک حفظ می‌شود.

\begin{figure}[H]
\centering
\includegraphics[height=6.2cm]{\detokenize{figures/shadow_game.png}}
\caption{تطبیق سه تصویر سایه‌ای با گزینه‌های تصویری در سطح سوم}
\label{Fig:UserShadowMatch}
\end{figure}

\subsection{بازی صورت و بدن}

بازی صورت و بدن کودک را در یک اجرای سه‌مرحله‌ای همراهی می‌کند. در هر مرحله، کودک قطعه‌های مشخص را می‌کشد و در محل مناسب روی همان شخصیت رها می‌کند. قطعهٔ درست در جای خود باقی می‌ماند و پاسخ نادرست بازخورد می‌گیرد؛ گزینه‌ها در طول اجرای یک شخصیت از همان مجموعه باقی می‌مانند. اجرای مرحله‌ها غیرتجمعی است: هر مرحله گروه تازه‌ای از مفاهیم را تمرین می‌کند و به بازچیدن قطعه‌های مرحله‌های پیشین نیاز ندارد.

مرحلهٔ نخست شامل چشم چپ، چشم راست و بینی است. در مرحلهٔ دوم، کودک ابروی چپ، ابروی راست و دهان را کامل می‌کند. مرحلهٔ سوم به دست چپ، دست راست، پای چپ و پای راست می‌پردازد. بنابراین هر مرحله به‌ترتیب سه، سه و چهار قطعه دارد. ستاره‌ها تکمیل مرحله‌های اجرا را نشان می‌دهند و بازنشانی دستی فقط قطعه‌های فعالیت جاری را پاک می‌کند. شخصیت در جریان هر سه مرحله ثابت می‌ماند. پس از جشن پایان مرحلهٔ سوم، اجرای تازه‌ای با شخصیت دیگری آغاز می‌شود؛ شخصیت‌های کامل فعلی برای این مسیر شامل شماره‌های ۰۳ تا ۰۷ هستند.

\begin{figure}[H]
\centering
\includegraphics[height=8cm]{\detokenize{figures/face_game.png}}
\caption{کامل‌کردن دست‌ها و پاهای شخصیت در مرحلهٔ سوم بازی صورت و بدن}
\label{Fig:UserFaceBody}
\end{figure}

\section{مدیریت محتوای شخصی}

والد می‌تواند از بخش مدیریت واژه‌ها، واژه‌های شخصی را همراه با تصویر و صدای ضبط‌شده ثبت و بعداً بازبینی یا حذف کند. برای افزودن واژه، نام فارسی، تصویر و فایل صوتی لازم‌اند؛ صدای واژه با میکروفن دستگاه ضبط می‌شود. در آغاز این فرایند، والد دسته‌ای را از فهرست جاری انتخاب می‌کند و در همان صفحه امکان ساخت دستهٔ تازه نیز فراهم است. ساخت دسته به نام فارسی نیاز دارد؛ تصویر و نام انگلیسی اختیاری‌اند.

دسته‌های پایه در فهرست محتوای برنامه مشترک‌اند و دستهٔ ساخته‌شده در این رابط خصوصیِ یک والد نیست. واژه‌های شخصی در مدیریت والد نمایش داده می‌شوند، اما در نسخهٔ جاری به‌طور خودکار به جریان عادی آموزش یا آزمون کودک افزوده نمی‌شوند. بنابراین افزودن واژه، به‌تنهایی محتوای برنامهٔ درسی را تغییر نمی‌دهد. والد باید هنگام انتخاب تصویر و ضبط صدا بر محتوا نظارت کند و از ثبت اطلاعات خصوصی غیرضروری بپرهیزد.

\begin{figure}[H]
\centering
\includegraphics[height=8cm]{\detokenize{figures/add_word_section1_categories.png}}
\caption{انتخاب دسته‌بندی در فرایند افزودن واژهٔ شخصی}
\label{Fig:UserAddWordCategorySelection}
\end{figure}

\begin{figure}[H]
\centering
\begin{minipage}{0.47\textwidth}
\centering
\includegraphics[height=6.3cm]{\detokenize{figures/add_word_section2_words.png}}
\\[-1mm]\small (الف) فهرست واژه‌های دستهٔ انتخاب‌شده
\end{minipage}
\hfill
\begin{minipage}{0.47\textwidth}
\centering
\includegraphics[height=6.3cm]{\detokenize{figures/add_word_section3_add_word.png}}
\\[-1mm]\small (ب) فرم تصویر و ضبط صوت واژه
\end{minipage}
\caption{مرور واژه‌های دسته و افزودن واژهٔ شخصی}
\label{Fig:UserWordManagement}
\end{figure}

\section{تنظیم زمان و حالت کودک}

در تنظیمات والد می‌توان مدت جلسهٔ آموزش را از میان یک، پنج، ده، پانزده یا سی دقیقه انتخاب کرد. این مقدار هنگام ورود به صفحهٔ آموزش به‌عنوان زمان همان جلسه به‌کار می‌رود و محدودیت زمانی یکپارچه برای همهٔ بازی‌ها نیست.

در Android، والد واردحساب‌شده می‌تواند از کنترل حالت کودک درخواست پین صفحه را به سیستم‌عامل بدهد. اگر پین در تنظیمات دستگاه فعال نباشد یا درخواست پذیرفته نشود، راهنمای درون‌برنامه‌ای مسیرهای متناسب با برخی برندها را نشان می‌دهد و وضعیت ناموفق را گزارش می‌کند. برداشتن پین از کنترل برنامه با تأیید دوبارهٔ رمز حساب والد انجام می‌شود؛ دسترسی به زبانهٔ تنظیمات نیز جداگانه به تأیید مجدد رمز نیاز دارد. این قابلیت از screen pinning خود Android استفاده می‌کند و برنامه سازوکار قفل سیستم‌عامل را پیاده‌سازی نمی‌کند. پین خودکار در iOS پشتیبانی نمی‌شود و راهنمای برنامه استفاده از قابلیت \lr{Guided Access} خود سیستم‌عامل را پیشنهاد می‌کند.

\begin{figure}[H]
\centering
\begin{minipage}{0.47\textwidth}
\centering
\includegraphics[height=7.2cm]{\detokenize{figures/app_pinning_main.png}}
\\[-1mm]\small (الف) راهنمای حالت کودک
\end{minipage}
\hfill
\begin{minipage}{0.47\textwidth}
\centering
\includegraphics[height=7.2cm]{\detokenize{figures/app_pinning_error_not_pinned.png}}
\\[-1mm]\small (ب) وضعیت پس از لغو درخواست Android
\end{minipage}
\caption{راهنمای درون‌برنامه‌ای و بازخورد در صورت فعال‌نشدن screen pinning}
\label{Fig:UserKidsModeGuide}
\end{figure}

\begin{figure}[H]
\centering
\includegraphics[height=8cm]{\detokenize{figures/app_pinning_pinned.png}}
\caption{صفحهٔ بازی هنگام فعال‌بودن screen pinning سیستم Android}
\label{Fig:UserKidsModePinned}
\end{figure}

حالت کودک یک کمک فنی برای محدودکردن خروج از برنامه است و جای نظارت والد را نمی‌گیرد. در Android، نام گزینه‌ها و مسیر فعال‌سازی ممکن است با سازنده و نسخهٔ سیستم‌عامل تفاوت داشته باشد.

\section{رفع مشکلات رایج}

\begin{itemize}
\item \textbf{تصویر یا صوت بارگذاری نمی‌شود:} اتصال اینترنت را بررسی کنید و پس از چند لحظه صفحه را دوباره باز کنید.
\item \textbf{صدای واژه پخش نمی‌شود:} صدای رسانه و حالت بی‌صدای دستگاه را بررسی کنید و پخش را دوباره امتحان نمایید.
\item \textbf{ضبط صدا یا انتخاب تصویر در بخش والد کار نمی‌کند:} مجوز میکروفن یا دسترسی به تصویرها را در تنظیمات دستگاه بررسی کنید.
\item \textbf{پیشرفت کودک دیده نمی‌شود:} از ورود به حساب والد و انتخاب نمایهٔ درست مطمئن شوید؛ همهٔ فعالیت‌ها پیشرفت پایدار ندارند.
\item \textbf{دسته یا واژه پیدا نمی‌شود:} اتصال شبکه و فهرست محتوای جاری را بررسی کنید. واژه‌های شخصی به‌طور خودکار در آموزش یا آزمون استاندارد ظاهر نمی‌شوند.
\item \textbf{پین برنامه در Android فعال نمی‌شود:} پین صفحه را در تنظیمات دستگاه فعال و راهنمای نمایش‌داده‌شده در برنامه را دنبال کنید.
\end{itemize}

\section{توصیه‌های استفادهٔ مناسب}

نوواژه برای ارائهٔ تمرین کوتاه و ایجاد فرصت گفت‌وگو طراحی شده است. همراهی والد، توجه به علاقه و خستگی کودک و پیوند دادن تصویرها به اشیا و موقعیت‌های واقعی، استفاده از برنامه را در بستر تعامل روزمره قرار می‌دهد. به‌ویژه برای کودک خردسال، زمان صفحه‌نمایش نباید جای بازی حرکتی، استراحت و تعامل پاسخ‌گوی والد را بگیرد؛ فعالیت را کوتاه نگه دارید و در صورت بی‌علاقگی کودک به پایان برسانید.

امتیاز یا شاخص‌های برنامه برای تشخیص رشد یا اختلال گفتار و زبان به‌کار نمی‌روند. اگر دربارهٔ شنوایی، درک زبان، تولید واژه یا از دست‌رفتن مهارت‌های کودک نگرانی وجود دارد، با پزشک یا آسیب‌شناس گفتار و زبان مشورت کنید.
